\documentclass[12pt]{article}
\usepackage[utf8]{inputenc}
\usepackage[T1]{fontenc}
\usepackage[margin=1in]{geometry}
\usepackage{setspace}
\usepackage{xcolor}
\usepackage[colorlinks=true, urlcolor=blue!60!black, linkcolor=blue!60!black]{hyperref}
\setlength{\parindent}{0em}
\setlength{\parskip}{0.6em}
\renewcommand{\baselinestretch}{1.0}
\begin{document}
\begin{center}
{\LARGE \textit{Buck v. Bell}, 274 U.S. 200 (1927)\par}
\vspace{0.5em}
{\large Case Brief --- Two-Page Edition\par}
\end{center}

\section{Memory Jogger}
The Court upheld Virginia's compulsory eugenic sterilization law, likening forced salpingectomy to compulsory vaccination (\href{https://supreme.justia.com/cases/federal/us/274/200/}{274 U.S. 200}).

\section{Detailed Case Facts}
Virginia's 1924 statute authorized sterilizing institutionalized ``mental defectives'' after a petition, notice to a guardian, a hearing, and two levels of appeal. Carrie Buck, eighteen, committed to the State Colony, was recorded as a ``feeble minded white woman'' with a ``feeble minded mother'' and child. Later research showed she had been raped and committed to hide the scandal, and neither she nor her daughter was intellectually disabled. Her counsel, a former colony board member, called no witnesses.

\section{Procedural History}
The colony board ordered sterilization; the Amherst County Circuit Court and the Virginia Supreme Court of Appeals affirmed; the U.S. Supreme Court affirmed in 1927.

\section{Judicial Votes}
Affirmed, 8--1. Holmes wrote for Taft, Van Devanter, McReynolds, Brandeis, Sutherland, Sanford, and Stone; Butler dissented without opinion.

\section{Holding}
Compulsory sterilization of institutional inmates adjudicated, after careful procedures, to carry hereditary insanity or imbecility is within the police power and satisfies due process; confining it to inmates does not deny equal protection, because a state may act incrementally.

\section{Analysis of Opinions}
\textbf{Majority:} The statute's procedures meant ``the rights of the patient are most carefully considered.'' Deferring to eugenic premises, Holmes reasoned that a state that can conscript its ``best citizens'' may demand ``these lesser sacrifices,'' and relied on \textit{Jacobson v. Massachusetts}. Equal protection was ``the usual last resort of constitutional arguments.'' \textbf{Butler:} dissented silently, leaving no counter-framework.

\section{Examples: Future Applications}
\textbf{Same-Side:} a Virginia-style vasectomy law for psychiatric patients with full hearings and appeals. \textbf{Opposite-Side:} a superintendent's summary sterilization with no notice or hearing fails \textit{Buck}'s procedural floor. \textbf{Fence-Sitter:} mandated reversible contraception, after hearings, for community-dwelling welfare recipients with hereditary disabilities.

\section{Critique}
\textit{Buck} belongs to the anti-canon: the vaccination analogy is a false equivalence, the record was a sham, and the equal protection reasoning was a rubber stamp. It fueled over 60,000 sterilizations and has never been formally overruled, though \textit{Skinner v. Oklahoma} gutted its rationale.

\section{Key Quotations}
\begin{quote}
``The principle that sustains compulsory vaccination is broad enough to cover cutting the Fallopian tubes. Jacobson v. Massachusetts, 197 U. S. 11. Three generations of imbeciles are enough.''
\end{quote}
\noindent --- Justice Holmes, Opinion of the Court, 274 U.S. at 207.

\end{document}
